%% Teaser (two-column strip). PLACEHOLDER panels; final content described in the caption and in README.md.
%% Planned panels, one ARKitScenes frame, clustered layout, k = 200, n = 20, UniDepthV2:
%%  (a) RGB with structured anchors (dots) and probes (crosses); (b) per-pixel miss map of A-local on the
%%  structured anchors; (c) miss map of A-dual; (d) per-image coverage histogram over test images, both methods,
%%  with the 0.9 target line.
\newcommand{\teaserpanel}[1]{%
  \begin{tikzpicture}
    \draw[gray!60, fill=gray!8] (0,0) rectangle (0.235\linewidth,0.15\linewidth);
    \node[align=center, font=\scriptsize\sffamily, text width=0.21\linewidth] at (0.1175\linewidth,0.075\linewidth) {#1};
  \end{tikzpicture}}
\begin{strip}
\centering
\teaserpanel{(a) RGB + structured anchors ($\bullet$) + random probes ($\times$)}\hfill
\teaserpanel{(b) \mLocal reusing structured anchors: misses (red) far from the cluster}\hfill
\teaserpanel{(c) \mDual: same anchors predict, probes calibrate}\hfill
\teaserpanel{(d) per-image coverage across test images; target 0.9}
\captionsetup{hypcap=false}
\captionof{figure}{\textbf{Reusing structured depth measurements for calibration fails silently; separating their roles does not.}
A frozen monocular model plus sparse metric anchors placed by the sensor (here one cluster).
Calibrating within the image on the same anchors (b) yields tight intervals that miss most of the frame;
spending a few measurements as uniformly random probes that never enter the prediction (c) restores the
per-image guarantee for any layout (d). \pending{placeholder panels: render from \texttt{experiments/qualitative.py}}}
\captionsetup{hypcap=true}
\label{fig:teaser}%
\end{strip}
